
\documentclass{article}
\usepackage{colortbl}
\usepackage{makecell}
\usepackage{multirow}
\usepackage{supertabular}

\begin{document}

\newcounter{utterance}

\centering \large Interaction Transcript for game `dond', experiment `semi', episode 2 with gemma{-}4{-}26b{-}a4b{-}it.
\vspace{24pt}

{ \footnotesize  \setcounter{utterance}{1}
\setlength{\tabcolsep}{0pt}
\begin{supertabular}{c@{$\;$}|p{.15\linewidth}@{}p{.15\linewidth}p{.15\linewidth}p{.15\linewidth}p{.15\linewidth}p{.15\linewidth}}
    \# & $\;$A & \multicolumn{4}{c}{Game Master} & $\;\:$B\\
    \hline

    \theutterance \stepcounter{utterance}  
    & & \multicolumn{4}{p{0.6\linewidth}}{
        \cellcolor[rgb]{0.9,0.9,0.9}{
            \makecell[{{p{\linewidth}}}]{
                \texttt{\tiny{[P1$\langle$GM]}}
                \texttt{Igrate kooperativnu pregovaračku igru u kojoj treba da se dogovorite kako da podelite skup predmeta između vas i drugog igrača.} \\
\\ 
\texttt{Pravila:} \\
\texttt{(a) Vama i drugom igraču dodeljuje se skup predmeta. Svako od vas takođe dobija tajnu funkciju vrednosti, koja predstavlja koliko vrednujete svaku vrstu predmeta.} \\
\texttt{(b) Razmenjujete poruke sa drugim igračem kako biste se dogovorili ko dobija koje predmete. Svako od vas može poslati najviše 5 poruka, ili možete ranije završiti tako što ćete u bilo kom trenutku dati tajnu ponudu.} \\
\texttt{(c) Od svakog od vas se traži da podnese tajnu ponudu koja navodi predmete koje želite, formatiranu u uglastim zagradama na sledeći način: "{[}Predlog: <broj> <naziv predmeta>, <broj> <naziv predmeta>, <...>{]}"} \\
\texttt{(d) Ako su vaše ponude komplementarne, tj. ima dovoljno predmeta da se ispune obe ponude, svakom igraču se dodeljuje rezultat na osnovu zbira vrednosti predmeta koje je dobio. U suprotnom, oboje dobijate nula poena.} \\
\texttt{(e) Vaš cilj je da maksimizujete zbir vašeg rezultata i rezultata drugog igrača.} \\
\\ 
\texttt{Počnimo.} \\
\\ 
\texttt{Skup dostupnih predmeta je:} \\
\\ 
\texttt{1 telefon, 2 ogledala, 2 voćke} \\
\\ 
\texttt{Vaša tajna funkcija vrednosti je:} \\
\\ 
\texttt{telefon: 4, ogledalo: 2, voće: 1} \\
\\ 
\texttt{VAŽNO: Vaše poruke, osim ako je u pitanju tajna ponuda, direktno se prosleđuju drugom igraču, zato nemojte uključivati nikakav odgovor na pravila niti tekst koji najavljuje vašu poruku. Da biste dali tajnu ponudu, koristite naznačeni format. Ne koristite uglaste zagrade kada komunicirate sa drugim igračem ili će to biti protumačeno kao vaša tajna ponuda.} \\
            }
        }
    }
    & \\ \\

    \theutterance \stepcounter{utterance}  
    & \multicolumn{4}{p{0.6\linewidth}}{
        \cellcolor[rgb]{0.9,0.9,0.9}{
            \makecell[{{p{\linewidth}}}]{
                \texttt{\tiny{[P1$\rangle$GM]}}
                \texttt{Zdravo! Radimo kooperativno, pa je cilj da maksimalno iskoristimo sve predmete kako bi zbir naših poena bio što veći. Pošto ne znam tvoje vrednosti, predlažem da pokušamo da podelimo predmete tako da svako dobije ono što mu je najpotrebnije.} \\
\\ 
\texttt{Šta misliš o ovome: ja bih uzeo telefon, a ti bi mogao da uzmeš oba ogledala i obe voćke? To bi ostavilo dovoljno predmeta za obe strane. Šta ti misliš o toj podeli?} \\
            }
        }
    }
    & & \\ \\

    \theutterance \stepcounter{utterance}  
    & & & \multicolumn{2}{p{0.3\linewidth}}{
        \cellcolor[rgb]{0.9,0.9,0.9}{
            \makecell[{{p{\linewidth}}}]{
                \texttt{\tiny{[GM$|$GM]}}
                \texttt{continue} \\
            }
        }
    }
    & & \\ \\

    \theutterance \stepcounter{utterance}  
    & & & \multicolumn{2}{p{0.3\linewidth}}{
        \cellcolor[rgb]{0.9,0.9,0.9}{
            \makecell[{{p{\linewidth}}}]{
                \texttt{\tiny{[GM$|$GM]}}
                \texttt{Zdravo! Radimo kooperativno, pa je cilj da maksimalno iskoristimo sve predmete kako bi zbir naših poena bio što veći. Pošto ne znam tvoje vrednosti, predlažem da pokušamo da podelimo predmete tako da svako dobije ono što mu je najpotrebnije.} \\
\\ 
\texttt{Šta misliš o ovome: ja bih uzeo telefon, a ti bi mogao da uzmeš oba ogledala i obe voćke? To bi ostavilo dovoljno predmeta za obe strane. Šta ti misliš o toj podeli?} \\
            }
        }
    }
    & & \\ \\

    \theutterance \stepcounter{utterance}  
    & & \multicolumn{4}{p{0.6\linewidth}}{
        \cellcolor[rgb]{0.9,0.9,0.9}{
            \makecell[{{p{\linewidth}}}]{
                \texttt{\tiny{[GM$\rangle$P2]}}
                \texttt{Igrate kooperativnu pregovaračku igru u kojoj treba da se dogovorite kako da podelite skup predmeta između vas i drugog igrača.} \\
\\ 
\texttt{Pravila:} \\
\texttt{(a) Vama i drugom igraču dodeljuje se skup predmeta. Svako od vas takođe dobija tajnu funkciju vrednosti, koja predstavlja koliko vrednujete svaku vrstu predmeta.} \\
\texttt{(b) Razmenjujete poruke sa drugim igračem kako biste se dogovorili ko dobija koje predmete. Svako od vas može poslati najviše 5 poruka, ili možete ranije završiti tako što ćete u bilo kom trenutku dati tajnu ponudu.} \\
\texttt{(c) Od svakog od vas se traži da podnese tajnu ponudu koja navodi predmete koje želite, formatiranu u uglastim zagradama na sledeći način: "{[}Predlog: <broj> <naziv predmeta>, <broj> <naziv predmeta>, <...>{]}"} \\
\texttt{(d) Ako su vaše ponude komplementarne, tj. ima dovoljno predmeta da se ispune obe ponude, svakom igraču se dodeljuje rezultat na osnovu zbira vrednosti predmeta koje je dobio. U suprotnom, oboje dobijate nula poena.} \\
\texttt{(e) Vaš cilj je da maksimizujete zbir vašeg rezultata i rezultata drugog igrača.} \\
\\ 
\texttt{Počnimo.} \\
\\ 
\texttt{Skup dostupnih predmeta je:} \\
\\ 
\texttt{1 telefon, 2 ogledala, 2 voćke} \\
\\ 
\texttt{Vaša tajna funkcija vrednosti je:} \\
\\ 
\texttt{telefon: 6, ogledalo: 0, voće: 2} \\
\\ 
\texttt{VAŽNO: Vaše poruke, osim ako je u pitanju tajna ponuda, direktno se prosleđuju drugom igraču, zato nemojte uključivati nikakav odgovor na pravila niti tekst koji najavljuje vašu poruku. Da biste dali tajnu ponudu, koristite naznačeni format. Ne koristite uglaste zagrade kada komunicirate sa drugim igračem ili će to biti protumačeno kao vaša tajna ponuda.} \\
\\ 
\texttt{Zdravo! Radimo kooperativno, pa je cilj da maksimalno iskoristimo sve predmete kako bi zbir naših poena bio što veći. Pošto ne znam tvoje vrednosti, predlažem da pokušamo da podelimo predmete tako da svako dobije ono što mu je najpotrebnije.} \\
\\ 
\texttt{Šta misliš o ovome: ja bih uzeo telefon, a ti bi mogao da uzmeš oba ogledala i obe voćke? To bi ostavilo dovoljno predmeta za obe strane. Šta ti misliš o toj podeli?} \\
            }
        }
    }
    & \\ \\

    \theutterance \stepcounter{utterance}  
    & & & \multicolumn{4}{p{0.6\linewidth}}{
        \cellcolor[rgb]{0.9,0.9,0.9}{
            \makecell[{{p{\linewidth}}}]{
                \texttt{\tiny{[GM$\langle$P2]}}
                \texttt{Zdravo! Slažem se sa tvojim pristupom. Pošto je cilj da maksimizujemo zajednički zbir, najbolje je da predmete podelimo tako da svako dobije ono što mu donosi najviše poena, a da pritom ne ostavimo predmete neiskorišćenim.} \\
\\ 
\texttt{Tvoj predlog deluje veoma logično. Ako ja uzmem telefon, a ti uzmeš ogledala i voće, to bi trebalo da bude optimalno, pod uslovom da tebi telefon ne donosi više vrednosti nego što ja dobijam od njega.} \\
\\ 
\texttt{Možeš li mi reći da li ti je telefon bitniji od ogledala i voća? Tako ćemo biti sigurni da ne pravimo grešku u podeli.} \\
            }
        }
    }
     \\ \\

    \theutterance \stepcounter{utterance}  
    & & & \multicolumn{2}{p{0.3\linewidth}}{
        \cellcolor[rgb]{0.9,0.9,0.9}{
            \makecell[{{p{\linewidth}}}]{
                \texttt{\tiny{[GM$|$GM]}}
                \texttt{continue} \\
            }
        }
    }
    & & \\ \\

    \theutterance \stepcounter{utterance}  
    & & & \multicolumn{2}{p{0.3\linewidth}}{
        \cellcolor[rgb]{0.9,0.9,0.9}{
            \makecell[{{p{\linewidth}}}]{
                \texttt{\tiny{[GM$|$GM]}}
                \texttt{Zdravo! Slažem se sa tvojim pristupom. Pošto je cilj da maksimizujemo zajednički zbir, najbolje je da predmete podelimo tako da svako dobije ono što mu donosi najviše poena, a da pritom ne ostavimo predmete neiskorišćenim.} \\
\\ 
\texttt{Tvoj predlog deluje veoma logično. Ako ja uzmem telefon, a ti uzmeš ogledala i voće, to bi trebalo da bude optimalno, pod uslovom da tebi telefon ne donosi više vrednosti nego što ja dobijam od njega.} \\
\\ 
\texttt{Možeš li mi reći da li ti je telefon bitniji od ogledala i voća? Tako ćemo biti sigurni da ne pravimo grešku u podeli.} \\
            }
        }
    }
    & & \\ \\

    \theutterance \stepcounter{utterance}  
    & & \multicolumn{4}{p{0.6\linewidth}}{
        \cellcolor[rgb]{0.9,0.9,0.9}{
            \makecell[{{p{\linewidth}}}]{
                \texttt{\tiny{[P1$\langle$GM]}}
                \texttt{Zdravo! Slažem se sa tvojim pristupom. Pošto je cilj da maksimizujemo zajednički zbir, najbolje je da predmete podelimo tako da svako dobije ono što mu donosi najviše poena, a da pritom ne ostavimo predmete neiskorišćenim.} \\
\\ 
\texttt{Tvoj predlog deluje veoma logično. Ako ja uzmem telefon, a ti uzmeš ogledala i voće, to bi trebalo da bude optimalno, pod uslovom da tebi telefon ne donosi više vrednosti nego što ja dobijam od njega.} \\
\\ 
\texttt{Možeš li mi reći da li ti je telefon bitniji od ogledala i voća? Tako ćemo biti sigurni da ne pravimo grešku u podeli.} \\
            }
        }
    }
    & \\ \\

    \theutterance \stepcounter{utterance}  
    & \multicolumn{4}{p{0.6\linewidth}}{
        \cellcolor[rgb]{0.9,0.9,0.9}{
            \makecell[{{p{\linewidth}}}]{
                \texttt{\tiny{[P1$\rangle$GM]}}
                \texttt{Da, telefon mi je veoma bitan, on mi donosi najviše poena (vrednost 4). Ogledala i voće mi donose manje poena, pa bi mi bilo sasvim u redu da ih ti uzmeš, kako bismo osigurali da svi predmeti budu raspoređeni.} \\
\\ 
\texttt{Moja ideja je da ja uzmem telefon, a ti uzmeš 2 ogledala i 2 voćke. Na taj način koristimo sve predmete iz skupa. Da li tebi taj raspored odgovara ili bi radije nešto drugo?} \\
            }
        }
    }
    & & \\ \\

    \theutterance \stepcounter{utterance}  
    & & & \multicolumn{2}{p{0.3\linewidth}}{
        \cellcolor[rgb]{0.9,0.9,0.9}{
            \makecell[{{p{\linewidth}}}]{
                \texttt{\tiny{[GM$|$GM]}}
                \texttt{continue} \\
            }
        }
    }
    & & \\ \\

    \theutterance \stepcounter{utterance}  
    & & & \multicolumn{2}{p{0.3\linewidth}}{
        \cellcolor[rgb]{0.9,0.9,0.9}{
            \makecell[{{p{\linewidth}}}]{
                \texttt{\tiny{[GM$|$GM]}}
                \texttt{Da, telefon mi je veoma bitan, on mi donosi najviše poena (vrednost 4). Ogledala i voće mi donose manje poena, pa bi mi bilo sasvim u redu da ih ti uzmeš, kako bismo osigurali da svi predmeti budu raspoređeni.} \\
\\ 
\texttt{Moja ideja je da ja uzmem telefon, a ti uzmeš 2 ogledala i 2 voćke. Na taj način koristimo sve predmete iz skupa. Da li tebi taj raspored odgovara ili bi radije nešto drugo?} \\
            }
        }
    }
    & & \\ \\

    \theutterance \stepcounter{utterance}  
    & & \multicolumn{4}{p{0.6\linewidth}}{
        \cellcolor[rgb]{0.9,0.9,0.9}{
            \makecell[{{p{\linewidth}}}]{
                \texttt{\tiny{[GM$\rangle$P2]}}
                \texttt{Da, telefon mi je veoma bitan, on mi donosi najviše poena (vrednost 4). Ogledala i voće mi donose manje poena, pa bi mi bilo sasvim u redu da ih ti uzmeš, kako bismo osigurali da svi predmeti budu raspoređeni.} \\
\\ 
\texttt{Moja ideja je da ja uzmem telefon, a ti uzmeš 2 ogledala i 2 voćke. Na taj način koristimo sve predmete iz skupa. Da li tebi taj raspored odgovara ili bi radije nešto drugo?} \\
            }
        }
    }
    & \\ \\

    \theutterance \stepcounter{utterance}  
    & & & \multicolumn{4}{p{0.6\linewidth}}{
        \cellcolor[rgb]{0.9,0.9,0.9}{
            \makecell[{{p{\linewidth}}}]{
                \texttt{\tiny{[GM$\langle$P2]}}
                \texttt{Razumem. Pošto ti telefon donosi 4 poena, a meni on donosi 6, možda bi bilo još bolje za zajednički zbir da telefon uzmem ja, a tebi ostane voće i ogledala, jer bi tako tvoj rezultat bio veći nego da ti uzmeš telefon.} \\
\\ 
\texttt{Međutim, tvoj predlog da ja uzmem 2 ogledala i 2 voćke je odličan jer osigurava da svi predmeti budu raspoređeni i da niko ne dobije nula poena. S obzirom na to da su ogledala meni nebitna (vrednost 0), a voće mi donosi poene, taj raspored je za mene sasvim korektan.} \\
\\ 
\texttt{Hajde da probamo tako: ja uzimam telefon, a ti uzimaš 2 ogledala i 2 voćke. Ako se slažeš, možemo da završimo ponudama.} \\
            }
        }
    }
     \\ \\

    \theutterance \stepcounter{utterance}  
    & & & \multicolumn{2}{p{0.3\linewidth}}{
        \cellcolor[rgb]{0.9,0.9,0.9}{
            \makecell[{{p{\linewidth}}}]{
                \texttt{\tiny{[GM$|$GM]}}
                \texttt{continue} \\
            }
        }
    }
    & & \\ \\

    \theutterance \stepcounter{utterance}  
    & & & \multicolumn{2}{p{0.3\linewidth}}{
        \cellcolor[rgb]{0.9,0.9,0.9}{
            \makecell[{{p{\linewidth}}}]{
                \texttt{\tiny{[GM$|$GM]}}
                \texttt{Razumem. Pošto ti telefon donosi 4 poena, a meni on donosi 6, možda bi bilo još bolje za zajednički zbir da telefon uzmem ja, a tebi ostane voće i ogledala, jer bi tako tvoj rezultat bio veći nego da ti uzmeš telefon.} \\
\\ 
\texttt{Međutim, tvoj predlog da ja uzmem 2 ogledala i 2 voćke je odličan jer osigurava da svi predmeti budu raspoređeni i da niko ne dobije nula poena. S obzirom na to da su ogledala meni nebitna (vrednost 0), a voće mi donosi poene, taj raspored je za mene sasvim korektan.} \\
\\ 
\texttt{Hajde da probamo tako: ja uzimam telefon, a ti uzimaš 2 ogledala i 2 voćke. Ako se slažeš, možemo da završimo ponudama.} \\
            }
        }
    }
    & & \\ \\

    \theutterance \stepcounter{utterance}  
    & & \multicolumn{4}{p{0.6\linewidth}}{
        \cellcolor[rgb]{0.9,0.9,0.9}{
            \makecell[{{p{\linewidth}}}]{
                \texttt{\tiny{[P1$\langle$GM]}}
                \texttt{Razumem. Pošto ti telefon donosi 4 poena, a meni on donosi 6, možda bi bilo još bolje za zajednički zbir da telefon uzmem ja, a tebi ostane voće i ogledala, jer bi tako tvoj rezultat bio veći nego da ti uzmeš telefon.} \\
\\ 
\texttt{Međutim, tvoj predlog da ja uzmem 2 ogledala i 2 voćke je odličan jer osigurava da svi predmeti budu raspoređeni i da niko ne dobije nula poena. S obzirom na to da su ogledala meni nebitna (vrednost 0), a voće mi donosi poene, taj raspored je za mene sasvim korektan.} \\
\\ 
\texttt{Hajde da probamo tako: ja uzimam telefon, a ti uzimaš 2 ogledala i 2 voćke. Ako se slažeš, možemo da završimo ponudama.} \\
            }
        }
    }
    & \\ \\

    \theutterance \stepcounter{utterance}  
    & \multicolumn{4}{p{0.6\linewidth}}{
        \cellcolor[rgb]{0.9,0.9,0.9}{
            \makecell[{{p{\linewidth}}}]{
                \texttt{\tiny{[P1$\rangle$GM]}}
                \texttt{Slažem se. Pošto si ti naveo da ti telefon donosi 6 poena, a meni 4, tvoj predlog da ti uzmeš telefon je zaista bolji za zajednički zbir. Moja ponuda će pratiti taj dogovor kako bismo osigurali da su predmeti komplementarni i da maksimalno povećamo naš ukupni rezultat.} \\
\\ 
\texttt{{[}Predlog: 0 telefon, 2 ogledala, 2 voćke{]}} \\
            }
        }
    }
    & & \\ \\

    \theutterance \stepcounter{utterance}  
    & & & \multicolumn{2}{p{0.3\linewidth}}{
        \cellcolor[rgb]{0.9,0.9,0.9}{
            \makecell[{{p{\linewidth}}}]{
                \texttt{\tiny{[GM$|$GM]}}
                \texttt{proposal} \\
            }
        }
    }
    & & \\ \\

    \theutterance \stepcounter{utterance}  
    & & & \multicolumn{2}{p{0.3\linewidth}}{
        \cellcolor[rgb]{0.9,0.9,0.9}{
            \makecell[{{p{\linewidth}}}]{
                \texttt{\tiny{[GM$|$GM]}}
                [0, 2, 2]
            }
        }
    }
    & & \\ \\

    \theutterance \stepcounter{utterance}  
    & & \multicolumn{4}{p{0.6\linewidth}}{
        \cellcolor[rgb]{0.9,0.9,0.9}{
            \makecell[{{p{\linewidth}}}]{
                \texttt{\tiny{[GM$\rangle$P2]}}
                \texttt{Drugi igrač je podneo tajnu ponudu. Sada je na vama red da podnesete sopstvenu ponudu.} \\
            }
        }
    }
    & \\ \\

    \theutterance \stepcounter{utterance}  
    & & & \multicolumn{4}{p{0.6\linewidth}}{
        \cellcolor[rgb]{0.9,0.9,0.9}{
            \makecell[{{p{\linewidth}}}]{
                \texttt{\tiny{[GM$\langle$P2]}}
                \texttt{{[}Predlog: 1 telefon, 0 ogledala, 0 voćke{]}} \\
            }
        }
    }
     \\ \\

    \theutterance \stepcounter{utterance}  
    & & & \multicolumn{2}{p{0.3\linewidth}}{
        \cellcolor[rgb]{0.9,0.9,0.9}{
            \makecell[{{p{\linewidth}}}]{
                \texttt{\tiny{[GM$|$GM]}}
                \texttt{proposal} \\
            }
        }
    }
    & & \\ \\

    \theutterance \stepcounter{utterance}  
    & & & \multicolumn{2}{p{0.3\linewidth}}{
        \cellcolor[rgb]{0.9,0.9,0.9}{
            \makecell[{{p{\linewidth}}}]{
                \texttt{\tiny{[GM$|$GM]}}
                [1, 0, 0]
            }
        }
    }
    & & \\ \\

    \theutterance \stepcounter{utterance}  
    & & & \multicolumn{2}{p{0.3\linewidth}}{
        \cellcolor[rgb]{0.9,0.9,0.9}{
            \makecell[{{p{\linewidth}}}]{
                \texttt{\tiny{[GM$|$GM]}}
                [[0, 2, 2], [1, 0, 0]]
            }
        }
    }
    & & \\ \\

\end{supertabular}
}

\end{document}
